\session{15}{11/21}{2限}{4}{総括：AI時代の経営情報}

\goals{
  \item 講義全体（4つのモジュール）の要点を、つながりを持って振り返る。
  \item 経営情報という分野のこれからと、人が担い続ける役割を自分の言葉で述べられる。
  \item 受講後に学びを続ける方法を具体的に持つ。
}

\guidesec{解説}

\topic{講義全体の振り返り}
本講義は、情報を経営資源として捉え、生成AIの活用を軸に、企業活動と情報技術の関係を学んできました。4つのモジュールの要点をつなげて振り返ります。
\par\noindent\begin{gtabh}{colspec={Q[l,wd=30mm] X[l] X[l]}}
  モジュール & 要点 & AIとの関わり \\
  01 経営と情報の基礎（第1〜3回） & 情報は第4の経営資源。意思決定は情報活動・設計活動・選択活動で進み、人の合理性には限界がある。生成AIは次のトークンを予測する確率的な道具で、正しさは保証されない。 & AIは代替案の生成と情報の整理で意思決定を支えるが、ハルシネーションを前提に検証が必要。 \\
  02 企業と情報システム（第4〜7回） & 企業活動はバリューチェーンで整理でき、ERP・CRM・SCM が支える。自動化は RPA から生成AI、AIエージェントへ広がり、人の仕事は「決める・確かめる・責任を持つ」に寄る。 & AIは SaaS として既存システムに組み込まれている。導入の成否は目的の明確さ、データ、体制、変革管理で決まる。 \\
  03 データと意思決定（第8〜11回） & データは集め、整え、可視化し、分析して判断に使う。相関は因果ではなく、予測は前提と限界を持つ。人の判断にはバイアスがある。 & AIは前処理・可視化・予測を速く行うが、問いの設定と解釈、人の関与の設計、説明責任は人が担う。 \\
  04 デジタル時代の経営（第12〜15回） & DX はビジネスモデルの変革。AI投資はコスト・効果・リスクで評価し、段階的に進める。AIにはバイアス・プライバシー・著作権・セキュリティのリスクがあり、ガバナンスが必要。 & AI活用提案は、課題と目的から始め、人の役割とリスクへの備えを含み、AIの寄与と自分の判断を区別して示す。 \\
\end{gtabh}
全体を貫くのは AI Policy の「使う・検証する・区別する」です。これは AI の使い方の規則であると同時に、情報を経営に活かす姿勢そのものです。

\topic{経営情報という分野のこれから}
生成AIの進化は速く、本講義で扱った技術の具体的な姿は数年で変わるでしょう。いくつかの動向を挙げます。
\begin{itemize}
  \item \textbf{AIエージェントの普及}　人が指示する「道具」から、目標を与えられて自律的に動く「担当者」に近づく。権限、監視、責任の設計（第5・11・13回）の重要性が増す。
  \item \textbf{汎用人工知能（AGI）をめぐる議論}　特定分野に限らず幅広い知的作業を人間並みにこなす AI を指す概念で、その定義や実現の時期は研究者の間でも見解が分かれる。経営の立場では、実現を前提にするのでも否定するのでもなく、能力の変化を継続的に評価する姿勢が必要。
  \item \textbf{規制と社会の要請}　AIの利用に対する法規制と、公平性・透明性への社会の要請は強まる。ガバナンスは「守り」ではなく、信頼を得て AI を使うための条件になる。
  \item \textbf{人の役割}　技術が変わっても、目的を定め、価値に基づいて判断し、結果に責任を持つことは人が担う。情報の質を見抜く力、問いを立てる力、説明する力は、むしろ価値が高まる。
\end{itemize}
経営情報という分野は、技術の変化を追いかけるだけでなく、「情報と技術を、経営の目的のためにどう使い、どう統制するか」という問いを持ち続ける分野です。
この問いは、技術が変わっても変わりません。

\topic{受講後に続けてほしい学び方}
\begin{itemize}
  \item \textbf{基本概念を軸にする}　新しい技術やサービスに出会ったら、情報（DIKW）、意思決定（3段階）、業務プロセス、データ品質、リスクとガバナンスの概念で整理する。概念は技術より長持ちする。
  \item \textbf{検証の習慣を続ける}　AIの出力の固有名詞・数値・文献を一次資料で確かめる。この習慣は、AI以外の情報（ニュース、営業資料、会議の発言）にも効く。
  \item \textbf{使い続け、記録する}　新しいツールを試し、何に使えて何に使えないかを自分の言葉で記録する。AI活用記録の形式は、職場でも使える。
  \item \textbf{一次情報に触れる}　官公庁の統計と公式文書、企業の公式資料、学術論文。AI利用ガイドの復習教材（先行研究の探し方、文献の批判的な読み方）は、この練習になる。
  \item \textbf{人と議論する}　AIとの対話は便利だが、人との議論でしか得られない視点がある。相互評価で経験したように、他者の視点は自分の提案を強くする。
\end{itemize}

\topic{キーワード}
\par\noindent\begin{keywords}{}
  使う・検証する・区別する & 本講義の AI Policy。情報を経営に活かす姿勢そのもの。 \\
  AIエージェント & 目標に向けて自律的に動く AI。権限・監視・責任の設計が重要になる。 \\
  汎用人工知能（AGI） & 幅広い知的作業を人間並みにこなす AI の概念。定義と実現時期は議論の対象。 \\
  人が担う役割 & 目的の設定、価値判断、結果への責任。技術が変わっても変わらない。 \\
  一次情報 & 統計、公式文書、論文など、AIの出力を検証するための独立した情報源。 \\
\end{keywords}

\guidesec{演習課題}

\exercise{1}{AI演習：講義全体をAIと共に振り返り、学んだことを自分の言葉でまとめる}
\instr{本ガイドの講義予定（はじめに）と、自分の演習の記録やメモを用意し、次のプロンプトでAIと対話しながら振り返ってください。AIの質問に答える形で進め、最後に「学んだこと」「考えが変わったこと」「これから続けること」を各200字程度で、自分の言葉で書いてください。AIの出力をそのまま使わないこと。}
\begin{promptbox}[振り返りの対話]
私は大学の「経営情報」という授業を15回受けました。授業は4つのモジュール（経営と情報の基礎、企業と情報システム、データと意思決定、デジタル時代の経営）で構成され、毎回生成AIを使う演習がありました。あなたは振り返りを手伝うコーチです。私に一度に1つずつ質問をして、「最も印象に残ったこと」「受講前と考えが変わったこと」「まだ分かっていないこと」「これから続けること」を引き出してください。私の答えを要約せず、掘り下げる質問を続けてください。
\end{promptbox}

\exercise{2}{概念で整理する}
\instr{最近のニュースから、企業のAI活用に関する記事を1つ選び、「情報（DIKW）」「意思決定の段階」「業務プロセスと自動化」「データ品質」「リスクとガバナンス」の5つの概念のうち3つ以上を使って、その記事の内容を整理してください。記事に書かれていないが確かめたいことも挙げます。}

\exercise{3}{学習計画を立てる}
\instr{受講後の半年間について、「続けること（習慣）」「試すこと（新しいツールや方法）」「読むこと（一次情報や文献）」を各2つ、いつ・どのように行うかを含めて書いてください。}

\guidesec{模範解答}

\begin{answer}{演習1}
例：
\begin{itemize}
  \item 学んだこと：生成AIは「検索」ではなく「生成」をしている道具であり、正しさは保証されないこと。だから出力を使う前に固有名詞と数値を一次資料で確かめる必要があること。
        AIを経営に使うには、技術より先に、解決したい課題と効果の測り方を決めなければならないこと。人の判断にもバイアスがあり、AIの推奨とどう向き合うかは設計の問題だということ。
  \item 考えが変わったこと：受講前は「AIを使えば仕事が速くなる」としか考えていなかったが、速くなった分、何を確かめ、何に責任を持つかを考える必要があると分かった。
        また、「AIを使うのはずるい」という感覚があったが、使い方を記録し明示すれば、AIの利用は評価される能力だと考えるようになった。
  \item これから続けること：AIの出力の固有名詞と数値を一次資料で確かめる習慣。レポートや就職活動の書類でAI活用記録を書くこと。興味のある企業の統合報告書を読んで、AI戦略を4観点で評価してみること。
\end{itemize}
評価の観点：AIの出力の要約ではなく、自分の経験（どの演習で何に気づいたか）に結びついた記述になっているか。
\end{answer}

\begin{answer}{演習2}
例：ある金融機関がコールセンターに生成AIを導入し、応対時間を2割短縮したという記事
\begin{itemize}
  \item 業務プロセスと自動化：応対中の回答候補の提示と応対後の要約が対象で、顧客との会話そのものは人が担う。生成AI向きの非定型業務に、人の確認を組み込んだ設計（Human-in-the-loop）。
  \item データ品質：過去の応対記録と FAQ が学習・検索の基盤。記事では FAQ の整備に半年かかったとあり、前処理の重要性を示している。
  \item リスクとガバナンス：顧客の個人情報を扱うため、入力の匿名化と学習に使わない契約が必要。記事には誤回答への対応が書かれていない。
  \item 確かめたいこと：応対時間2割短縮の測定方法（導入前後の比較か、導入店舗と非導入店舗の比較か）、顧客満足度の変化、誤った案内が起きたときの手順。
\end{itemize}
\end{answer}

\begin{answer}{演習3}
例：
\begin{itemize}
  \item 続けること：(1) AIの出力を使うときは固有名詞と数値を一次資料で確かめ、確認した資料をメモに残す（毎回）。(2) 月に1回、その月に試したAIの使い方と、うまくいかなかったことを記録する（月末）。
  \item 試すこと：(1) 表計算のデータを AI にグラフ化させ、第9回の観点で誤解を招く点を探す（11月中に1回）。(2) AIエージェント型のツールで、調べものから報告書の下書きまでを任せ、どこで人の確認が必要になったかを記録する（12月）。
  \item 読むこと：(1) 興味のある企業2社の統合報告書の AI・DX の部分を読み、第12回の演習の観点で比べる（1月）。(2) AI利用ガイドの復習教材「先行研究の探し方」に沿って、関心のある問いについて論文を5件探して読む（2〜3月）。
\end{itemize}
\end{answer}

\guidesec{出席確認クイズの解説}
講義サイトの出席確認ページ（第15回）の5問です。
% 出席確認クイズ 第15回　総括：AI時代の経営情報
%   attendance/attendance15.html から生成（解説は本ガイド用に加筆）。

\quizq{1}{本講義で学んだ「経営情報」の要点として適切なものはどれでしょうか?}%
  {AIにすべて任せれば経営判断は不要}%
  {AIは使わない方がよい}%
  {\correct{情報を経営資源として捉え、AIを検証しながら意思決定に活かす}}%
  {情報は経営とは無関係}%
  {正解は (c)。情報を経営資源として捉え、AIの出力を検証しながら意思決定に活かすことが本講義の中心テーマでした。(a) AI に任せきりにすることも、(b) AI を使わないことも、(d) 情報を経営と無関係とみなすことも、本講義の立場ではありません。}

\quizq{2}{生成AIの出力を確かめる方法として適切なものはどれでしょうか?}%
  {同じAIにもう一度聞くだけ}%
  {長い回答なら信じる}%
  {確かめる必要はない}%
  {\correct{一次情報や公式資料など、別の情報源で裏付けを取る}}%
  {正解は (d)。AIの出力は、一次情報や公式資料など、AIとは独立した情報源で裏付けを取って確かめます。(a) 同じAIに再質問しても同じ誤りを繰り返すことがあり、(b) 回答の長さは正確さの根拠になりません。}

\quizq{3}{AI時代に人が担い続けると考えられる役割として適切なものはどれでしょうか?}%
  {\correct{目的の設定、価値判断、結果への責任}}%
  {計算の代行}%
  {定型作業の反復だけ}%
  {何も担わない}%
  {正解は (a)。目的を定め、価値に基づいて判断し、結果に責任を持つことは、AI ではなく人が担う役割です。(b)(c) は機械に置き換えられやすい作業で、(d) は本講義の立場に反します。}

\quizq{4}{「汎用人工知能(AGI)」の説明として適切なものはどれでしょうか?}%
  {\correct{特定の課題に限らず、幅広い知的作業を人間並みにこなすAI(の概念)}}%
  {表計算ソフトの機能}%
  {画像だけを扱うAI}%
  {すでに一般に普及している家電}%
  {正解は (a)。AGI（Artificial General Intelligence、汎用人工知能）は、特定分野に限らず幅広い知的作業を人間並みにこなせる AI を指す概念で、その定義や実現時期は現在も研究・議論の対象です。(b)(c)(d) は AGI ではありません。}

\quizq{5}{継続的に学ぶ姿勢として適切なものはどれでしょうか?}%
  {流行を無批判に追いかける}%
  {\correct{新しいAIの動向を追いつつ、基本概念と検証の習慣を持ち続ける}}%
  {AIの話題は避ける}%
  {一度学べば更新は不要}%
  {正解は (b)。AIは変化が速いため、動向を追いながらも、基本概念（情報、意思決定、データ、ガバナンス）と出力を検証する習慣を持ち続けることが大切です。(a) 無批判な追随と (d) 更新の放棄はどちらも望ましくなく、(c) 話題を避けることは学びの機会を失います。}

